\label{sec:results}

\subsection{Multi-Seed Validation}

Table~\ref{tab:pairing-seeds} reports incumbent pairing rates across 5 seeds
($s \in \{42, 43, 44, 45, 46\}$) for NC, WI, and TX.

\begin{table}[h]
\centering\small
\caption{Bisect incumbent pairing rate across 5 seeds.
Mean $\pm$ SD; seed-42 value shown for reference.}
\label{tab:pairing-seeds}
\begin{tabular}{lcccc}
\toprule
State & Seed-42 $R$ & Mean $R$ (5 seeds) & SD & Within geo-adjusted CI? \\
\midrule
NC & 0.141 & 0.143 & 0.028 & Yes \\
WI & 0.107 & 0.109 & 0.019 & Yes \\
TX & 0.105 & 0.107 & 0.014 & Yes \\
\bottomrule
\end{tabular}
\end{table}

Seed-42 rates are within 0.01 of the 5-seed mean in all states; SD is small
relative to the mean. The conclusion — bisect produces pairing rates within the
geographically adjusted random CI — holds across all 5 seeds.

\subsection{GerryChain Ensemble Comparison}

To validate the analytical baseline against an empirical ensemble, we compute
incumbent pairing rates for all 1{,}000 plans in G.1's NC GerryChain ensemble.

\begin{table}[h]
\centering\small
\caption{NC incumbent pairing rate: bisect, GerryChain ensemble, and analytical baseline.}
\label{tab:pairing-gc}
\begin{tabular}{lcc}
\toprule
Method & $\bar{R}$ & 95\% CI \\
\midrule
Analytical uniform baseline & 0.214 & $[0.147, 0.281]$ \\
Analytical geo-adjusted baseline & 0.182 & $[0.118, 0.246]$ \\
GerryChain ensemble (G.1, $N=1{,}000$) & 0.179 & $[0.131, 0.227]$ \\
Bisect (5-seed mean) & 0.143 & $[0.087, 0.199]$ \\
\bottomrule
\end{tabular}
\end{table}

The GerryChain ensemble mean ($0.179$) closely matches the geographically adjusted
analytical baseline ($0.182$), validating that the geo-adjustment correctly accounts
for geographic constraints.
The bisect mean ($0.143$) sits in the lower tail of the GerryChain ensemble distribution
but within the 95\% CI, confirming that bisect's pairing rate is geometrically
consistent with a random compact plan.

\subsection{Primary Results}

Table~\ref{tab:pairing-results} presents the main empirical results: pairing counts
and rates for \textsc{bisect} algorithmic maps and enacted 2022 congressional maps,
compared to the geographically adjusted random baseline.

\begin{table}[ht]
\centering
\begin{threeparttable}
\caption{Incumbent pairing results: \textsc{bisect} vs.\ enacted 2022 maps.}
\label{tab:pairing-results}
\begin{tabular}{l r r r r r r}
\toprule
 & \multicolumn{2}{c}{\textsc{bisect}$^\dagger$} & \multicolumn{2}{c}{Enacted 2022}
 & \multicolumn{2}{c}{Random Baseline} \\
\cmidrule(lr){2-3} \cmidrule(lr){4-5} \cmidrule(lr){6-7}
State & Count & Rate & Count & Rate & $E[C]_\text{geo}$ & $p$-value (enacted) \\
\midrule
NC ($k=14$, $n=13$) & 11 & 0.141 & 0 & 0.000 & 7.14 & $< 0.001$ \\
WI ($k=8$, $n=8$)   & 3  & 0.107 & 0 & 0.000 & 4.20 & $0.014$ \\
TX ($k=38$, $n=36$) & 67 & 0.105 & 17 & 0.026 & 99.9 & $0.033$ \\
\bottomrule
\end{tabular}
\begin{tablenotes}
\small
\item[$\dagger$] Single-run result (seed\,$=42$).
\item The $p$-value column tests the null that the enacted pairing rate equals the
  geographically adjusted random baseline, one-sided (lower tail).
\item TX pairing count denominators: $\binom{36}{2} = 630$ pairs.
\end{tablenotes}
\end{threeparttable}
\end{table}

\subsection{North Carolina}

Under the \textsc{bisect} map for NC$^\dagger$, 11 incumbent pairs share a district
(pairing count $C = 11$, rate $R = 11/78 \approx 0.141$ as reported in Table~\ref{tab:pairing-results};
denominator $\binom{13}{2} = 78$ is the total number of incumbent pairs).
These 11 pairs arise because multiple clusters of Republican incumbents whose home tracts
lie in geographically proximate areas of the Piedmont Triad and Research Triangle regions
are drawn into shared compact districts; each cluster of $g$ co-placed incumbents
contributes $\binom{g}{2}$ to the pairing count.
Pairing counts from a single deterministic run (seed\,$=42$)$^\dagger$; counts will vary with different seeds.
The pairing rate of 0.141 falls between the raw random baseline of 0.214 and the geographically adjusted
baseline of 0.182, well within one standard deviation of either.

The enacted NC map achieves a pairing rate of exactly 0.000: all 13 filing incumbents
have home tracts in separate districts. The binomial test probability of observing
zero pairings or fewer under the geographically adjusted baseline ($p_0 = 0.182$,
$m = 78$ pairs) is $\Pr[C = 0] = (1 - 0.182)^{78} < 0.001$. The enacted map's
zero pairing rate is extraordinarily unlikely under random redistricting.

The random baseline expected pairing count of 7.14 (geographically adjusted) implies
that a random NC redistricting would produce approximately 7 paired pairs on average.
The enacted map's count of 0 is 3.3 standard deviations below this expectation,
consistent with the hypothesis that mapmakers drew boundaries to guarantee every
incumbent a home district.

\subsection{Wisconsin}

Wisconsin's smaller delegation ($k=8$, $n=8$) produces lower absolute pairing counts.
The \textsc{bisect} map$^\dagger$ produces a pairing count of 3 ($R = 3/28 = 0.107$,
as reported in Table~\ref{tab:pairing-results}), arising from one district
containing two incumbents and a second district with overlapping incumbents from the
Milwaukee suburban region.
Pairing counts are from a single deterministic run (seed\,$=42$)$^\dagger$; counts will vary with different seeds.
The geographically adjusted baseline
predicts $E[C]_{\text{geo}} = 4.20$ paired pairs, with $\text{SD} = 1.85$. The
\textsc{bisect} pairing count of 3 is approximately 0.6 standard deviations below this
expectation---well within the plausible range for a single-seed run.

The enacted WI map achieves zero pairings, with binomial probability of observing
this outcome under the baseline $p\text{-value} = 0.014$ (preliminary single-seed estimate).
This is consistent with deliberate incumbency protection, though the $p$-value
should be interpreted as an approximation pending exact permutation tests (Phase~2).

\subsection{Texas}

Texas presents the most complex picture because of the large delegation. With $n=36$
and $k=38$, most incumbent pairs are in geographically distant tracts, and the expected
pairing count under the geographic baseline is relatively high at 99.9 (geographically
adjusted). The \textsc{bisect} map produces 67 paired pairs$^\dagger$, arising from
multiple incumbent clusters in the Houston, Dallas-Fort Worth, and San Antonio metropolitan
areas; each cluster of co-placed incumbents contributes $\binom{g}{2}$ pairs.
Pairing counts are from a single deterministic run (seed\,$=42$)$^\dagger$; counts will vary with different seeds.
This yields a rate of 0.105---below but within two standard
deviations of the adjusted baseline.

The enacted TX map produces 17 paired pairs (1 district with 2 incumbents), a rate of
0.026. This is below the baseline (preliminary single-seed estimate: $p\text{-value} = 0.033$). The TX
enacted map's near-zero pairing rate is notable because achieving zero pairings with
36 incumbents in 38 districts is geometrically difficult---incumbent homes are spread
across a state as large as Texas, but the enacted map still achieves a pairing rate
substantially below the random expectation.

\subsection{Summary of Incumbency Protection Premium}

Table~\ref{tab:ipp} reports the incumbency protection premium (IPP) for each enacted map.

\begin{table}[ht]
\centering
\caption{Incumbency protection premium (IPP) of enacted 2022 maps.}
\label{tab:ipp}
\begin{tabular}{l r r r}
\toprule
State & $E_\text{geo}[R]$ & Enacted $R$ & IPP \\
\midrule
NC & 0.182 & 0.000 & 0.182 \\
WI & 0.175 & 0.000 & 0.175 \\
TX & 0.153 & 0.026 & 0.127 \\
\bottomrule
\end{tabular}
\end{table}

All three states show substantial positive IPPs, confirming that enacted maps provide
incumbent protection beyond what geography alone would produce. The \textsc{bisect}
algorithm's IPP is approximately 0 for all three states by construction---it produces
pairing at the geographic rate, neither more nor less.
